\documentclass[10pt,a4paper]{amsart}
\usepackage{amsmath,amssymb,amsthm,graphicx}
\usepackage{verbatim}
\newenvironment{reference}{\comment}{\endcomment}
%\newenvironment{reference}{}{}
\newenvironment{slogan}{\comment}{\endcomment}
\newenvironment{history}{\comment}{\endcomment}

% For commutative diagrams we use Xy-pic
\usepackage[all]{xy}

% We use 2cell for 2-commutative diagrams.
\xyoption{2cell}
% Theorem environments.
%
\theoremstyle{plain}
\newtheorem{theorem}[subsection]{Theorem}
\newtheorem{proposition}[subsection]{Proposition}
\newtheorem{lemma}[subsection]{Lemma}

\theoremstyle{definition}
\newtheorem{definition}[subsection]{Definition}
\newtheorem{example}[subsection]{Example}
\newtheorem{exercise}[subsection]{Exercise}
\newtheorem{situation}[subsection]{Situation}

\theoremstyle{remark}
\newtheorem{remark}[subsection]{Remark}
\newtheorem{remarks}[subsection]{Remarks}

\numberwithin{equation}{subsection}

% Macros
%
\def\lim{\mathop{\mathrm{lim}}\nolimits}
\def\colim{\mathop{\mathrm{colim}}\nolimits}
\def\Spec{\mathop{\mathrm{Spec}}}
\def\Hom{\mathop{\mathrm{Hom}}\nolimits}
\def\Ext{\mathop{\mathrm{Ext}}\nolimits}
\def\SheafHom{\mathop{\mathcal{H}\!\mathit{om}}\nolimits}
\def\SheafExt{\mathop{\mathcal{E}\!\mathit{xt}}\nolimits}
\def\Sch{\mathit{Sch}}
\def\Mor{\mathop{\mathrm{Mor}}\nolimits}
\def\Ob{\mathop{\mathrm{Ob}}\nolimits}
\def\Sh{\mathop{\mathit{Sh}}\nolimits}
\def\NL{\mathop{N\!L}\nolimits}
\def\CH{\mathop{\mathrm{CH}}\nolimits}
\def\proetale{{pro\text{-}\acute{e}tale}}
\def\etale{{\acute{e}tale}}
\def\QCoh{\mathit{QCoh}}
\def\Ker{\mathop{\mathrm{Ker}}}
\def\Im{\mathop{\mathrm{Im}}}
\def\Coker{\mathop{\mathrm{Coker}}}
\def\Coim{\mathop{\mathrm{Coim}}}

% Boxtimes
%
\DeclareMathSymbol{\boxtimes}{\mathbin}{AMSa}{"02}

%
% Macros for moduli stacks/spaces
%
\def\QCohstack{\mathcal{QC}\!\mathit{oh}}
\def\Cohstack{\mathcal{C}\!\mathit{oh}}
\def\Spacesstack{\mathcal{S}\!\mathit{paces}}
\def\Quotfunctor{\mathrm{Quot}}
\def\Hilbfunctor{\mathrm{Hilb}}
\def\Curvesstack{\mathcal{C}\!\mathit{urves}}
\def\Polarizedstack{\mathcal{P}\!\mathit{olarized}}
\def\Complexesstack{\mathcal{C}\!\mathit{omplexes}}
% \Pic is the operator that assigns to X its picard group, usage \Pic(X)
% \Picardstack_{X/B} denotes the Picard stack of X over B
% \Picardfunctor_{X/B} denotes the Picard functor of X over B
\def\Pic{\mathop{\mathrm{Pic}}\nolimits}
\def\Picardstack{\mathcal{P}\!\mathit{ic}}
\def\Picardfunctor{\mathrm{Pic}}
\def\Deformationcategory{\mathcal{D}\!\mathit{ef}}
\usepackage{cleveref}
\crefname{originallemma}{original lemma}{original lemmas}
\Crefname{originallemma}{Original lemma}{Original lemmas}
\makeatletter
\newlabel{lemma-graded-module-fg}{{00K0}{source}{}{sourcepointer-1}{}}
\newlabel{lemma-graded-module-fg@cref}{{[originallemma][0][]00K0}{[1][0][]source}{}{}{}}
\makeatother
\crefname{originallemma}{original lemma}{original lemmas}
\Crefname{originallemma}{Original lemma}{Original lemmas}
\makeatletter
\newlabel{lemma-numerical-polynomial}{{00JZ}{source}{}{sourcepointer-2}{}}
\newlabel{lemma-numerical-polynomial@cref}{{[originallemma][0][]00JZ}{[1][0][]source}{}{}{}}
\makeatother
\title{The Hilbert function of a finite graded module is a numerical polynomial in degree one}
\author{The Stacks Project Authors}
\date{stacks-project:algebra.tex:local-snapshot-20261012; GNU FDL 1.2 or later}
\begin{document}
\maketitle
\noindent\textit{Editorial note.} All delivered mathematical text is the author's own, copied from the pinned Stacks Project chapter file algebra.tex; no statement or proof step was rewritten, reordered or removed.\par\medskip
\noindent\textit{Editorial note.} Preamble provenance, as required for this source: the frozen snapshot stages the chapter's own preamble file beside it (preamble.tex, 217 lines, 8116 bytes, the same shared Stacks-wide preamble the published chow and etale-cohomology cards use); the first line of the pinned chapter file is \textbackslash{}input\{preamble\}, so no preamble text lives inside the chapter itself. The card takes that file's lines 8-18 (the verbatim package and the three \textbackslash{}comment-based environments reference/slogan/history at lines 9-12, section 13.2, together with the Xy-pic setup) and its lines 154 to the runtime end of the file (the native theorem declarations and mathematical macros). Not one character was changed, no package was added, no macro was defined, and the excerpt serves only as preamble, never as a span coordinate.\par\medskip
\noindent\textit{Editorial note.} The same-chapter pointers are located at their \textbackslash{}label lines inside the pinned chapter file; the cross-chapter pointers are located in the project's own tag list staged in the snapshot (sources/stacks-chow/tags.txt), at the verbatim lines carrying their canonical tags.\par\medskip
\noindent\textit{Editorial note.} The counted statement and proof refer to same-chapter statements by label; they are recorded as navigation stubs pointed at their label lines.\par\medskip
\noindent\textit{Editorial note.} The counted spans contain no cite, so the card carries no bibliography.\par\medskip
\noindent\textit{Allowed standard prerequisites (editorial).}\begin{itemize}
\item Finiteness of the graded pieces over the degree-zero ring. Invoked in the statement itself at line 14115. Makes the classes of the graded pieces elements of the Grothendieck group.
\item Numerical polynomials via the difference operator. Invoked at the end of the counted proof (line 14160). Converts the degree-by-degree difference computation into the polynomial conclusion.
\end{itemize}\medskip
\ifdefined\hypertarget\hypertarget{sourcepointer-1}{}\fi
\noindent\textit{Source reference.} Original lemma 00K0: sources/stacks-algebra/pinned-original-source.tex, source lines 14090--14090. Navigation pointer for \textbackslash{}ref\{lemma-graded-module-fg\} used inside the counted span (line 14115); the argument that consumes it is retained in the span.\par\medskip
\ifdefined\hypertarget\hypertarget{sourcepointer-2}{}\fi
\noindent\textit{Source reference.} Original lemma 00JZ: sources/stacks-algebra/pinned-original-source.tex, source lines 14070--14070. Navigation pointer for \textbackslash{}ref\{lemma-numerical-polynomial\} used inside the counted span (line 14160); the argument that consumes it is retained in the span.\par\medskip
\begin{proposition}
\label{proposition-graded-hilbert-polynomial}
Suppose that $S$ is a Noetherian graded ring
and $M$ a finite graded $S$-module. Consider the
function
$$
\mathbf{Z} \longrightarrow K'_0(S_0), \quad
n \longmapsto [M_n]
$$
see Lemma \ref{lemma-graded-module-fg}.
If $S_{+}$ is generated by elements of degree $1$,
then this function is a numerical polynomial.
\end{proposition}

\begin{proof}
We prove this by induction on the minimal number of
generators of $S_1$. If this number is $0$, then
$M_n = 0$ for all $n \gg 0$ and the result holds.
To prove the induction step, let $x\in S_1$
be one of a minimal set of generators, such that
the induction hypothesis applies to the
graded ring $S/(x)$.

\medskip\noindent
First we show the result holds if $x$ is nilpotent on $M$.
This we do by induction on the minimal integer $r$ such that
$x^r M  = 0$. If $r = 1$, then $M$ is a module over $S/xS$
and the result holds (by the other induction hypothesis).
If $r > 1$, then we can find a short exact sequence
$0 \to M' \to M \to M'' \to 0$ such that the integers
$r', r''$ are strictly smaller than $r$. Thus we know
the result for $M''$ and $M'$. Hence
we get the result for $M$ because of the relation
$
[M_d]  = [M'_d] + [M''_d]
$
in $K'_0(S_0)$.

\medskip\noindent
If $x$ is not nilpotent on $M$, let $M' \subset M$ be
the largest submodule on which $x$ is nilpotent.
Consider the exact sequence $0 \to M' \to M \to M/M' \to 0$
we see again it suffices to prove the result for $M/M'$. In other
words we may assume that multiplication by $x$ is injective.

\medskip\noindent
Let $\overline{M} = M/xM$. Note that the map $x : M \to M$
is {\it not} a map of graded $S$-modules, since it does
not map $M_d$ into $M_d$. Namely, for each $d$ we have the
following short exact sequence
$$
0 \to M_d \xrightarrow{x} M_{d + 1} \to \overline{M}_{d + 1} \to 0
$$
This proves that $[M_{d + 1}] - [M_d] = [\overline{M}_{d + 1}]$.
Hence we win by Lemma \ref{lemma-numerical-polynomial}.
\end{proof}

\end{document}
